
%%%%%%%%%%%%%%%%%%%%%%%%%%%%%%%%%%%%%%%%%%%%%%%%%%%%%%%%%%%%%
\subsubsection{具体分析}

问题一要求分析蔬菜各品类及单品销售量的分布规律及相互关系，属于描述性分析和统计推断类问题。本问采用Spearman秩相关系数进行品类间关联分析，以K-means++聚类分析单品销售模式，并通过经验分布拟合刻画销售量的分布形态。

首先对销售流水按品类和日期聚合为日销售量矩阵（1049天$\times$6品类），剔除退货（461条，0.05\%）和5个从未销售的单品。然后计算Spearman秩相关系数量化品类间关联，以轮廓系数确定最优聚类数对单品分组，最后对各品类日销售量拟合候选分布（正态、对数正态、Gamma）并以KS检验选择最优分布。

\subsubsection{模型准备}

在进行建模之前，需要对原始销售流水数据进行聚合和清洗。原始数据来源于附件2，包含878,503条交易记录，涵盖2020年7月1日至2023年6月30日。

经检查发现以下问题：退货461条（0.05\%），5个单品在全部三年中无任何销售记录，87个单品销售天数不足30天。针对上述问题，按以下步骤进行预处理：第一步，剔除退货记录和从未销售的单品；第二步，将交易数据按品类和日期聚合为日销售量；第三步，单品聚类分析仅纳入销售天数$\geq 30$的单品（159/246），以确保特征估计的可靠性。预处理后，得到1049天$\times$6品类的日销售量矩阵。

该问题本质上是一个探索性数据分析问题，需要在数据非正态的前提下选择合适的统计方法。全部分个品类均拒绝D'Agostino-Pearson正态性检验（$p < 10^{-75}$），因此Pearson相关系数假设不成立。Spearman秩相关系数作为非参数方法，对分布无要求且对异常值稳健，是更合适的选择。K-means++聚类通过优化初始质心选择可得到比标准K-means更稳定的分组结果。因此本文选择Spearman秩相关作为主要关联分析方法，K-means++作为聚类算法。

经上述预处理，得到1049天×6品类的日销售量矩阵和159个稳定单品的特征向量。接下来建立Spearman秩相关和K-means++聚类模型。
